\chapter{Mi firma digital}\label{cap:mi-firma-digital}

En \textbf{Mi firma digital} usted prepara la clave y el certificado con los
que firma documentos y contratos. El recorrido completo es:

\begin{enumerate}
  \item Generar su clave y su solicitud de certificado (CSR) en el navegador.
  \item Guardar la clave cifrada y comprobar la copia.
  \item Imprimir la hoja de huella y entregar la solicitud al responsable de
        la autoridad certificadora (CA) del despacho.
  \item Recibir su certificado y vincularlo a su cuenta.
\end{enumerate}

\disponible Administrador, Litigante y Cliente.

\importante{Su clave privada y su contraseña nunca salen de su equipo. Qadra
no puede recuperarlas: si las pierde, deberá generar una clave nueva. Los
certificados son de la CA interna del despacho y no sustituyen a la e.firma
del SAT.}

% ---------------------------------------------------------------------
\section{Generar la clave y la solicitud}\label{sec:generar-clave}

\paraque Crea su clave privada protegida con contraseña y la solicitud
pública (CSR) que necesita la CA para emitir su certificado.

\pasos
\begin{enumerate}
  \item En el menú, haga clic en \textbf{Mi firma digital}. Si es
        Administrador, quédese en la pestaña \textbf{Firma de documentos}.
  \item En \textbf{Generar clave y solicitud}, llene
        \textbf{Datos de la solicitud}: \textbf{Nombre común (CN)},
        \textbf{Organización (O)} y \textbf{País (C)} (por ejemplo,
        \texttt{MX}).
  \item Escriba la \textbf{Contraseña de la nueva clave} (al menos 12
        caracteres) y repítala en \textbf{Confirmar contraseña}.
  \item Haga clic en \textbf{Generar clave y solicitud}.
  \item Haga clic en \textbf{Descargar clave cifrada} y guarde el archivo
        \archivo{qadra-document-signing-key.encrypted.pem} en un lugar seguro.
  \item Para comprobar la copia, selecciónela en
        \textbf{Archivo de clave cifrada guardado}, escriba su contraseña en
        \textbf{Contraseña de la clave guardada} y haga clic en
        \textbf{Comprobar copia guardada}.
  \item Haga clic en \textbf{Descargar solicitud CSR}. Se descarga
        \archivo{qadra-document-signing-request.csr.pem}.
\end{enumerate}

\begin{figure}[H]
  \centering
  \capturapendiente{capturas/A-40-generar-clave.png}%
    {Sección Generar clave y solicitud con Datos de la solicitud y contraseña · Litigante}
  \caption{Generación de la clave y la solicitud.}
  \label{fig:manual-generar-clave}
\end{figure}

\resultado Aparece \enquote{La copia cifrada se puede abrir y corresponde a
esta solicitud.} y queda descargada su solicitud CSR.

\nota{El nombre que escribe es declarado: no acredita su identidad. Para el
Cliente, la organización es fija y corresponde a la CA emisora.}

% ---------------------------------------------------------------------
\section{Hoja de huella y entrega de la solicitud}\label{sec:hoja-huella}

\paraque Asegura que el certificado que emitirá la CA corresponde a la clave
que usted generó.

\pasos
\begin{enumerate}
  \item En \textbf{Huella de tu clave pública}, haga clic en
        \textbf{Imprimir hoja de huella}.
  \item Imprima la \textbf{Hoja de huella de firma documental}. Tiene dos
        ejemplares: \textbf{Ejemplar del Cliente} y
        \textbf{Ejemplar del despacho}.
  \item Firmen a mano ambos ejemplares (usted y el despacho). Anote el
        \textbf{Folio}; se pide al solicitar las firmas de un contrato.
  \item Entregue solo el archivo CSR al responsable de la CA.
  \item Lea la huella al responsable de la CA, en persona o por teléfono.
        Debe coincidir con la que él ve antes de emitir el certificado.
\end{enumerate}

\begin{figure}[H]
  \centering
  \capturapendiente{capturas/A-41-hoja-huella.png}%
    {Vista de impresión de la Hoja de huella de firma documental con ambos ejemplares (datos de prueba) · Litigante}
  \caption{Hoja de huella de firma documental.}
  \label{fig:manual-hoja-huella}
\end{figure}

\resultado La CA le entrega su certificado público (archivo \texttt{.pem}).

\importante{Comparta solo el CSR. Nunca entregue su clave privada ni su
contraseña, ni siquiera al responsable de la CA.}

% ---------------------------------------------------------------------
\section{Registrar el certificado y el vínculo}\label{sec:registrar-vinculo}

\paraque Vincula su certificado a su cuenta de Qadra. Sin vínculo no puede
firmar documentos ni contratos.

\pasos
\begin{enumerate}
  \item En \textbf{Mi firma digital}, vaya a
        \textbf{Registrar certificado documental}.
  \item Seleccione su certificado en \textbf{Certificado público PEM} y haga
        clic en \textbf{Preparar vínculo}.
  \item Revise los datos del certificado.
  \item En \textbf{Firmar en este dispositivo}, seleccione su
        \textbf{Clave privada cifrada}, escriba la
        \textbf{Contraseña de la clave} y haga clic en
        \textbf{Firmar localmente}.
  \item Marque \textbf{Confirmo vincular este certificado a mi cuenta}.
  \item Haga clic en \textbf{Registrar vínculo}.
\end{enumerate}

\begin{figure}[H]
  \centering
  \capturapendiente{capturas/A-42-registrar-vinculo.png}%
    {Registrar certificado documental con certificado cargado y Firmar en este dispositivo · Litigante}
  \caption{Registro del vínculo del certificado.}
  \label{fig:manual-registrar-vinculo}
\end{figure}

\resultado Aparece \textbf{Vínculo registrado}. Haga clic en
\textbf{Descargar recibo del vínculo} para conservar el recibo
(\archivo{certificate-binding-receipt.json}).

\nota{Si prefiere firmar con su propia herramienta, abra
\textbf{Usar una firma externa}, descargue la declaración con
\textbf{Descargar declaración del vínculo}, fírmela y cargue el resultado en
\textbf{Firma externa del vínculo}.}

% ---------------------------------------------------------------------
\section{Retirar el vínculo}\label{sec:retirar-vinculo}

\paraque Deja de usar un certificado, por ejemplo si perdió su clave o va a
renovarlo.

\importante{Retirar el vínculo impide firmar con ese certificado en adelante.
Las firmas anteriores conservan su evidencia.}

\pasos
\begin{enumerate}
  \item En \textbf{Vínculo registrado}, haga clic en \textbf{Retirar vínculo}.
  \item Marque \textbf{Confirmo retirar el vínculo para futuras firmas}.
  \item Haga clic en \textbf{Confirmar retiro}.
\end{enumerate}

\resultado Aparece \textbf{Vínculo retirado}. Para volver a firmar, genere
una clave nueva o registre un certificado renovado.

% ---------------------------------------------------------------------
\section{Certificado vencido o por vencer}\label{sec:certificado-vence}

Qadra le avisa en \textbf{Mi firma digital}:

\begin{itemize}
  \item \textbf{Por vencer:} \enquote{Tu certificado vence en menos de 30
        días. Renuévalo y vuelve a vincularlo; una solicitud abierta con el
        certificado anterior debe rehacerse.}
  \item \textbf{Vencido:} \enquote{Tu certificado venció. Genera una clave
        nueva o renueva el certificado y vuelve a vincularlo.}
\end{itemize}

En ambos casos, repita las \cref{sec:generar-clave,sec:hoja-huella,sec:registrar-vinculo}.

% ---------------------------------------------------------------------
\section{Certificado para ingresar (Administrador)}\label{sec:acceso-certificado}

\paraque Registra el certificado con el que el Administrador puede entrar a
Qadra (\cref{sec:ingresar-certificado}). Este certificado no firma
documentos.

\disponible Administrador.

\pasos
\begin{enumerate}
  \item En \textbf{Mi firma digital}, abra la pestaña
        \textbf{Acceso con certificado}.
  \item Genere su clave y su solicitud igual que en la
        \cref{sec:generar-clave}, y obtenga su certificado de la CA.
  \item Haga clic en \textbf{Registrar certificado}.
  \item Seleccione su \textbf{Certificado público PEM} y haga clic en
        \textbf{Preparar vínculo}.
  \item Firme en \textbf{Firmar en este dispositivo} con su clave y su
        contraseña, y haga clic en \textbf{Registrar vínculo}.
  \item En \textbf{Vínculo de mi cuenta}, haga clic en
        \textbf{Descargar recibo}. Ese archivo \texttt{.json} es el que
        seleccionará al ingresar con certificado.
\end{enumerate}

\begin{figure}[H]
  \centering
  \capturapendiente{capturas/A-43-acceso-certificado.png}%
    {Pestaña Acceso con certificado con Vínculo de mi cuenta y Descargar recibo · Administrador}
  \caption{Certificado para ingresar a Qadra.}
  \label{fig:manual-acceso-certificado}
\end{figure}

\resultado El vínculo aparece como \textbf{Sin retirar} y ya puede usar
\textbf{Ingresar con certificado}.

\nota{Para dejar de usarlo, haga clic en \textbf{Retirar vínculo} y luego en
\textbf{Confirmar retiro}. El retiro es definitivo para ese vínculo y cierra
las sesiones abiertas con él; su acceso con contraseña y segundo factor no
cambia.}
